% ===== PRÉAMBULE CANONIQUE — à coller VERBATIM en tête de CHAQUE .tex =====
\documentclass[11pt,a4paper]{article}
\usepackage[utf8]{inputenc}\usepackage[T1]{fontenc}\usepackage{lmodern}
\usepackage[french]{babel}
\usepackage{amsmath,amssymb,mathtools}\usepackage{icomma}
\usepackage{siunitx}\sisetup{locale=FR}  % \qty{}{}, \si{}, \ang{}
\usepackage{xcolor}\usepackage{colortbl}
\usepackage{tikz}\usepackage{pgfplots}\pgfplotsset{compat=1.18}
\usepackage[siunitx,europeanresistors]{circuitikz} % circuits (élec.)
\usepackage[most]{tcolorbox}\usepackage{enumitem}\usepackage{array,booktabs}
\usepackage[a4paper,margin=2.2cm]{geometry}
\usetikzlibrary{arrows.meta,calc,angles,quotes,positioning,patterns,%
  backgrounds,shapes.geometric,decorations.pathmorphing,decorations.markings,intersections}
\definecolor{primary}{RGB}{222,49,99}\definecolor{primarydark}{RGB}{150,22,55}
\definecolor{defcol}{RGB}{0,114,178}\definecolor{methcol}{RGB}{52,92,148}
\definecolor{excol}{RGB}{0,135,90}\definecolor{attncol}{RGB}{200,40,55}
\tcbset{boxbase/.style={enhanced,breakable,boxrule=0.9pt,arc=2.5pt,
  left=3mm,right=3mm,top=2.3mm,bottom=2.3mm,fonttitle=\bfseries,coltitle=white}}
% --- boîtes TUTORIEL (noms EXACTS, ne pas renommer) ---
\newtcolorbox{cle}[1][]{boxbase,colback=defcol!6,colframe=defcol,
  title={\textsc{À retenir}\ifx\relax#1\relax\relax\else\textmd{ : \itshape #1}\fi}}
\newtcolorbox{methode}[1][]{boxbase,colback=methcol!7,colframe=methcol,sharp corners,
  title={\textsc{Méthode}\ifx\relax#1\relax\relax\else\textmd{ --- \itshape #1}\fi}}
\newtcolorbox{exemplebox}[1][]{boxbase,colback=excol!5,colframe=excol,
  title={\textsc{Exemple}\ifx\relax#1\relax\relax\else\textmd{ : \itshape #1}\fi}}
\newtcolorbox{piege}[1][]{boxbase,colback=attncol!6,colframe=attncol,
  title={\textsc{Piège}\ifx\relax#1\relax\relax\else\textmd{ : \itshape #1}\fi}}
\newtcolorbox{remarque}[1][]{boxbase,colback=black!4,colframe=black!45,
  title={\textsc{Remarque}\ifx\relax#1\relax\relax\else\textmd{ : \itshape #1}\fi}}
% --- boîtes EXERCICE / CORRIGÉ (noms EXACTS, auto-comptées) ---
\newtcolorbox[auto counter]{exo}[1][]{boxbase,colback=primary!4,colframe=primary,
  title={\textsc{Exercice}~\thetcbcounter\ifx\relax#1\relax\relax\else\textmd{ --- \itshape #1}\fi}}
\newtcolorbox[auto counter]{corrige}[1][]{boxbase,colback=excol!4,colframe=excol,
  title={\textsc{Corrigé}~\thetcbcounter\ifx\relax#1\relax\relax\else\textmd{ --- \itshape #1}\fi}}
\newcommand{\vv}[1]{\vec{#1}}\newcommand{\ds}{\displaystyle}
\newcommand{\R}{\mathbb{R}}\newcommand{\N}{\mathbb{N}}\newcommand{\Z}{\mathbb{Z}}
% (un \usepackage{fancyhdr}+en-tête est possible ; il est ignoré côté web)
% =========================================================================
\begin{document}
\begin{corrige}[quelles barres sont en équilibre ?]
\begin{enumerate}[label=\alph*)]
  \item Une barre posée sur un pivot est en équilibre si son centre de gravité $G$ est \textbf{à la
        verticale du pivot} (ici, directement au-dessus). Sinon, le poids (appliqué en $G$) possède un
        bras de levier non nul par rapport au pivot : il crée un moment qui fait basculer la barre.
  \item \textbf{Barre 1} : $G$ est au-dessus du pivot $\Rightarrow$ \textbf{équilibre}. \textbf{Barre 2} :
        $G$ est à droite du pivot $\Rightarrow$ le poids la fait basculer \emph{vers la droite}.
        \textbf{Barre 3} : $G$ est à gauche $\Rightarrow$ elle bascule \emph{vers la gauche}.
\end{enumerate}
\end{corrige}

\begin{corrige}[centre de gravité d'une plaque en L]
\begin{enumerate}[label=\alph*)]
  \item $R_1$ : aire $A_1 = 6\times 2 = \qty{12}{cm^2}$, centre $G_1=(3\,;\,1)$.
        $R_2$ : aire $A_2 = 2\times 4 = \qty{8}{cm^2}$, centre $G_2=(1\,;\,4)$.
  \item On pondère par les aires :
        \[ x_G = \frac{A_1 x_1 + A_2 x_2}{A_1+A_2} = \frac{12\times 3 + 8\times 1}{20}
        = \frac{44}{20} = \qty{2.2}{cm}, \]
        \[ y_G = \frac{A_1 y_1 + A_2 y_2}{A_1+A_2} = \frac{12\times 1 + 8\times 4}{20}
        = \frac{44}{20} = \qty{2.2}{cm}. \]
        Le centre de gravité est en $G=(\qty{2.2}{cm}\,;\,\qty{2.2}{cm})$ — dans le creux du L, à la
        limite de la matière.
\end{enumerate}
\end{corrige}

\begin{corrige}[à partir de quelle inclinaison bascule-t-il ?]
\begin{enumerate}[label=\alph*)]
  \item Le point $G$ est au centre : il est à $b/2$ de l'arête basse (le long de la base) et à $h/2$
        au-dessus. Le bloc bascule quand la verticale de $G$ passe par cette arête, c'est-à-dire quand
        la droite (arête $\to G$) devient verticale. Cette droite fait avec la perpendiculaire à la base
        un angle $\arctan\!\big(\tfrac{b/2}{h/2}\big)=\arctan\!\big(\tfrac{b}{h}\big)$ ; à la bascule, la base
        est inclinée d'exactement cet angle, donc
        \[ \tan\theta_c = \frac{b}{h}. \]
  \item $\tan\theta_c = \dfrac{30}{60} = 0{,}5 \Rightarrow \theta_c = \arctan(0{,}5)\approx\ang{26.6}$.
  \item Élargir le bloc (augmenter $b$) \emph{augmente} $\theta_c$ : plus stable. Le rehausser
        (augmenter $h$) le \emph{diminue} : moins stable. \textbf{Mieux vaut l'élargir} — base large et
        $G$ bas.
\end{enumerate}
\end{corrige}

\begin{corrige}[placer une masse pour viser un centre de gravité]
\begin{enumerate}[label=\alph*)]
  \item $x_G = \dfrac{m_1\times 0 + m_2\times x}{m_1+m_2} = 40 \Rightarrow \dfrac{4x}{5}=40
        \Rightarrow 4x = 200 \Rightarrow x = \qty{50}{cm}$.
  \item \textbf{Oui.} Suspendue par son centre de gravité, la règle a son poids qui passe exactement par
        le point de suspension : le bras de levier est nul, aucun moment ne la fait tourner. Elle reste
        horizontale (équilibre indifférent).
\end{enumerate}
\end{corrige}

\end{document}
